% ============================================================================
% AOGS EXPERT BRIEFING — presenter's defense dossier for the poster session.
% Companion to aogs_study_explainer.pdf (plain language). THIS document is
% aimed at EXPERT questioners: every equation explained, figure-by-figure
% talking points, and a Q&A bank built from a simulated 5-reviewer panel
% (2026-07-12) whose attacks mirror what HFE-literate visitors will ask.
% Numbers via poster_numbers.tex (generated, never hand-typed) where macros
% exist; remaining values traced to the results archives on the dated audit.
% Build:  /Library/TeX/texbin/pdflatex aogs_expert_briefing.tex   (x2)
% ============================================================================
\documentclass[11pt]{article}
\usepackage[a4paper, margin=20mm]{geometry}
\usepackage[T1]{fontenc}
\usepackage{newtxtext,newtxmath}
\usepackage{graphicx}
\usepackage{booktabs}
\usepackage{xcolor}
\usepackage[font=small,labelfont=bf]{caption}
\usepackage{enumitem}
\usepackage{amsmath}
\usepackage{float}
\newcommand{\horAfif}{14.0}
\newcommand{\lossAfif}{1.16}
\newcommand{\svfAfif}{0.985}
\newcommand{\kdstarAfif}{4.60}
\newcommand{\horAsev}{10.1}
\newcommand{\lossAsev}{0.18}
\newcommand{\svfAsev}{0.991}
\newcommand{\kdstarAsev}{7.08}
\newcommand{\rmsehayAfif}{2.03}
\newcommand{\rmsecpAfif}{1.74}
\newcommand{\bestzcpAfif}{5/30}
\newcommand{\bestrhocpAfif}{1500/1700/1700}
\newcommand{\rmsecKAfif}{0.90}
\newcommand{\bestzcKAfif}{10/30}
\newcommand{\bestrhocKAfif}{1500/1700/1800}
\newcommand{\rmsehayAsev}{0.54}
\newcommand{\rmsecpAsev}{0.39}
\newcommand{\bestzcpAsev}{5/30}
\newcommand{\bestrhocpAsev}{1500/1700/1700}
\newcommand{\rmsecKAsev}{0.36}
\newcommand{\bestzcKAsev}{5/30}
\newcommand{\bestrhocKAsev}{1300/1500/1900}
\newcommand{\rmsehomAfif}{2.31}
\newcommand{\rmsehomAsev}{3.76}
\newcommand{\xferatAsev}{1.64}
\newcommand{\xferatAfif}{1.73}
\newcommand{\svfdtirAfif}{1.15}
\newcommand{\svfdtallAfif}{1.30}
\newcommand{\svfdtirAsev}{0.65}
\newcommand{\svfdtallAsev}{0.74}
\newcommand{\shadowcoolAfif}{2.22}
\newcommand{\shadowcoolAsev}{0.48}
\newcommand{\stabcKAfif}{13}
\newcommand{\bootcicKAfif}{0.26--1.08}
\newcommand{\stabcKAsev}{36}
\newcommand{\bootcicKAsev}{0.25--0.43}
\newcommand{\driftmax}{35}
\newcommand{\closurewin}{0.70}


\definecolor{teal}{HTML}{2A6478}
\definecolor{coral}{HTML}{B85B3A}
\definecolor{forest}{HTML}{3D6E4A}
\definecolor{char}{HTML}{2A2520}
\definecolor{gridc}{HTML}{E8E5E0}
\definecolor{dim}{HTML}{8A857E}
\color{char}

\graphicspath{{../figures/poster/}{../figures/diagrams/}{../figures/explainers/}}
\setlist[itemize]{leftmargin=1.4em, itemsep=1pt, topsep=2pt}
\setlist[description]{itemsep=2pt, topsep=2pt}

\newcommand{\shead}[1]{\medskip\noindent{\color{teal}\large\bfseries #1}\par\nobreak\smallskip}
% Q&A blocks
\newcommand{\Q}[1]{\par\medskip\noindent{\color{coral}\bfseries Q. #1}\par\nobreak\smallskip}
\newcommand{\A}{\noindent{\color{teal}\bfseries A.}\ }
% inline "say this" card
\newenvironment{saybox}%
{\par\smallskip\noindent\begin{minipage}{\linewidth}\color{char}%
 \colorbox{gridc}{\begin{minipage}{0.97\linewidth}\vspace{3pt}\small}%
{\vspace{3pt}\end{minipage}}\end{minipage}\par\smallskip}

\begin{document}
\thispagestyle{empty}

\begin{center}
{\LARGE\bfseries The AOGS Study, Expert Edition}\\[2pt]
{\large\color{teal}\bfseries Presenter's briefing: physics, figures, and the hard questions}\\[4pt]
{\small Ram\'on III P.\ Gregorio \quad$\bullet$\quad Institute of Science Tokyo}\\[-2pt]
{\footnotesize\color{dim} companion to the plain-language explainer; Q\&A built from a
 simulated five-reviewer panel (2026-07-12)}
\end{center}
\vspace{2pt}
{\color{teal}\rule{\linewidth}{0.9pt}}

% ── 1. The story ────────────────────────────────────────────────────────────
\shead{1\enspace The whole study in five sentences}
The only in-situ measurements of the lunar subsurface thermal state are the
Apollo 15 and 17 Heat Flow Experiment records. We take a certified 1-D
thermal-diffusion model of those sites and add two pieces of physics the
standard description lacks: \textbf{discrete density layers} (instead of the
smooth exponential of Hayne et al.\ 2017) and \textbf{topographic shadowing}
ray-traced from the LOLA digital elevation model. Sweeping 650 layer
configurations against the deep sensors shows the record prefers a surface
denser than Hayne's surface value and a shallow transition to consolidated
regolith --- and the fit improves not because density stores heat differently,
but because density \emph{implies conductivity} through the model's own
$K_c(\rho)$ relation. The retrieved deep densities (1800/1900 kg\,m$^{-3}$)
match the Apollo drive-core measurements (${\sim}$1825/1960 kg\,m$^{-3}$;
Grott et al.) in both magnitude and site ordering. Topography, meanwhile,
enters through two opposing channels of comparable size --- horizon blocking
cools the deep sensors (\shadowcoolAfif/\shadowcoolAsev\ K), terrain
re-radiation warms them ($+\svfdtallAfif$/$+\svfdtallAsev$ K upper bound) ---
so neither can be ignored on the assumption that the other cancels it.

\shead{2\enspace The claim ladder --- what we assert, at what strength}
Experts will test whether you know the difference between your strong and weak
claims. Hold this ladder in mind; never defend a rung higher than the evidence.
\begin{enumerate}[leftmargin=1.6em, itemsep=2pt]
  \item \textbf{Strong (defend to the end):} the deep-sensor record is
        reproduced to \rmsecKAfif/\rmsecKAsev\ K by a 3-layer column; the
        conductivity channel dominates the heat-capacity channel
        (${\sim}4$ K vs $<0.8$ K of leverage); the two topographic channels
        are individually $\gtrsim$1 K-scale and oppose.
  \item \textbf{Solid:} the improvement over baselines is bootstrap-robust
        (coupled best-RMSE 95\% CIs \bootcicKAfif/\bootcicKAsev\ K, entirely
        below homogeneous \rmsehomAfif/\rmsehomAsev\ K); retrieved deep
        densities agree with the drive cores.
  \item \textbf{Scoped (state the qualifier yourself):} against the
        \emph{field-standard} Hayne continuous model the gains are
        ${\sim}2.3\times$ (A15) and ${\sim}1.5\times$ (A17) --- the
        ${\sim}2.6\times/10\times$ figures are vs.\ homogeneous. At A17 the
        cp branch (\rmsecpAsev\ K) nearly ties coupled (\rmsecKAsev\ K).
  \item \textbf{Class-level only:} the A15 \emph{exact} best triple
        (\bestrhocKAfif) survives only \stabcKAfif\% of bootstrap draws; what
        recurs is the class --- denser surface, shallow consolidated
        transition. Never defend the individual triple.
  \item \textbf{Diagnostic / next step (concede freely):} terrain-radiation
        numbers are an upper bound ($T_{\rm terr}\!\approx\!T_s$) not yet
        folded into the sweep; the 16 ppd DEM under-resolves near-field
        terrain; $K_d$--deep-density degeneracy needs a joint retrieval.
\end{enumerate}

% ── 3. equations ────────────────────────────────────────────────────────────
\shead{3\enspace Every equation on the poster, and what to say about it}

\textbf{Heat equation.}\quad
$\rho(z)\,c_p(T)\,\partial_t T=\partial_z\!\big[K(T,z)\,\partial_z T\big]$.
1-D vertical conduction; no convection or vapor transport in dry regolith.
\emph{Say:} ``everything site-specific enters through the coefficients and
the boundary conditions.''

\textbf{Conductivity form (Hayne et al.\ 2017).}\quad
$K(T,z)=K_c(z)\,[1+\chi\,(T/350\,\mathrm{K})^3]$ with
$K_c(z)=K_d-(K_d-K_s)\,e^{-z/H}$; $\chi=2.7$, $K_s=0.74$, $H=6$ cm.
The bracket is the \emph{radiative} term --- photons hopping between grains,
hence the $T^3$; $K_c$ is the contact (phonon) part rising from a loose
surface to a dense deep value $K_d$.
\emph{Say:} ``we keep the standard form; we change what sets $K_c$.''

\textbf{The site-specific anchor.}\quad
$K_d^{\star}(\mathrm{A15})=\kdstarAfif$, $K_d^{\star}(\mathrm{A17})=\kdstarAsev$
mW\,m$^{-1}$K$^{-1}$ --- the deep conductivities certified by the parent
retrieval at each site. The \textbf{cp branch} holds $K_d=K_d^{\star}$ fixed
(density moves only $\rho c_p$). The \textbf{coupled branch} sets
$K_d=K_c(\rho_3)$ from the line through $(\rho_s{=}1100,\,K_s)$ and
$(\rho_d{=}1800,\,K_d^{\star})$ --- no new constants. Consequence: the A17
winner's $\rho_3=1900$ gives an \emph{effective} deep conductivity
$K_c(1900)=7.99$, not $7.08$.
\emph{Say:} ``the two branches are a controlled experiment on what density is
allowed to touch.''

\textbf{Boundary conditions.}\quad Surface:
$(1{-}A)\,S_{\rm sh}(t)=\varepsilon\sigma T_s^4+K\,\partial_z T|_0$
(absorbed sunlight balances thermal emission plus conduction). Base:
$K\,\partial_z T|_{\rm deep}=Q_b=21/16$ mW\,m$^{-2}$ (A15/A17; Langseth
et al.\ 1976 revised values).

\textbf{Shadowed insolation.}\quad
$S_{\rm sh}(t)=(1{-}A)\,S_0\cos\theta_s$ if $\theta_s>\theta_h(\varphi)$,
else $0$: the direct beam is deleted whenever the Sun sits below the
DEM-derived skyline $\theta_h$ at its current azimuth $\varphi$. Horizons are
built by marching rays outward in 90 azimuths, keeping the maximum elevation
angle with the curvature drop $r^2/2R$ subtracted.

\textbf{Flux-anchored closure.}\quad Below the rectification zone
($z_0=55$ cm) the periodic-equilibrium mean profile must satisfy
$d\langle T\rangle/dz=(Q_b-u_{\rm rect})/K$, where $u_{\rm rect}$ is the
rectified downward flux generated above $z_0$ by the nonlinearity (the
$T^3$ term responds more to hot than to cold extremes). This closure is what
makes one solve take ${\sim}0.6$ s (${\sim}20\times$ faster than brute-force
time-stepping to equilibrium); winners close the flux budget to
$\le\closurewin\%$ and agree to $\le\driftmax$ mK at $n{=}96$ grid points.
\emph{Say:} ``the sign is minus $u_{\rm rect}$ --- rectification steals part
of the basal flux from the mean gradient.''

\textbf{Sky-view factor.}\quad $\mathrm{SVF}=\overline{\cos^2\theta_h}$
(azimuth average): the fraction of sky visible from the site
(\svfAfif/\svfAsev). The complement $(1-\mathrm{SVF})$ is filled by terrain
that radiates thermal IR and reflects sunlight back at the surface --- the
warming channel.

\textbf{The score.}\quad
$\mathrm{RMSE}=\sqrt{\langle(\overline{T}(z_i)-T_i^{\rm obs})^2\rangle_i}$
over the deep sensors (0.8--2.4 m; shallower sensors excluded as
borestem-perturbed), with one equilibrium temperature per sensor from quiet
late-mission windows of the 1971--77 record (the 1975--77 portion restored by
Nagihara et al.\ 2018).

% ── 4. figures ─────────────────────────────────────────────────────────────
\shead{4\enspace Figure-by-figure: the one thing to point at}

\begin{description}
  \item[Fig 1 (schematic).] Point at the two numbered tags: layers (1) and
    shadowing (2) --- ``these two are the contribution; everything else is the
    certified baseline.'' Likely question: why three layers? --- coarsest
    structure that can hold a loose skin, a transition, and a consolidated
    base; boundaries and densities are swept, not assumed.
  \item[Fig 2 (pipeline).] Point at the teal-vs-grey legend. Walk the chain:
    horizons $\to$ layers $\to$ solver $\to$ mean profile $\to$ RMSE. The
    box to know cold is step 2 (the two branches and the anchor).
  \item[Fig 3 (densities + transfer).] Point at the staircase vs the dotted
    Hayne curve: ``the record wants more density up here.'' Then the matrix
    diagonal (\rmsecKAfif/\rmsecKAsev) vs off-diagonal
    (\xferatAsev/\xferatAfif): transfer degrades gracefully and still beats
    homogeneous. \emph{Concede if pushed:} at A17 the transplanted structure
    does not beat Hayne continuous (0.54 K) --- the transfer argument is
    strongest at A15 (1.73 K vs Hayne's 2.03 K).
  \item[Table 1 (RMSE ladder).] Read it left to right as ``models of
    increasing physics.'' Have the two honest ratios ready: vs.\ homogeneous
    ${\sim}2.6\times/10\times$; vs.\ Hayne ${\sim}2.3\times/1.5\times$.
  \item[Fig 4 (sensitivity).] Point at the width of the two box plots:
    ``same density values, two different doors --- through $c_p$ nothing
    happens; through $K$ everything happens.'' The maps below: best fits sit
    at $z_2=30$ cm at both sites, though A17 admits a tied deeper solution
    --- say that yourself before anyone reads it off the map.
  \item[Fig 5 (validation profiles).] Point at the slopes: the coupled
    branch \emph{tilts} the deep gradient onto the sensors; cp can only
    shift a fixed slope. That geometric statement is the mechanism in one
    sentence.
  \item[Fig 6 (bootstrap).] Point at the whole CI sitting left of the
    dashed/arrowed homogeneous value. If asked about the A15 cp interval
    (upper end 2.17 K, close to \rmsehomAfif\ K): the robust statement is
    about the coupled branch; the cp branch at A15 is marginal --- and
    that is itself evidence that heat capacity alone is not enough.
\end{description}

% ── 5. numbers card ────────────────────────────────────────────────────────
\shead{5\enspace Numbers card (certified values)}
\begin{center}\small
\begin{tabular}{@{}lcc@{}}
\toprule
 & \textbf{Apollo 15} & \textbf{Apollo 17}\\
\midrule
$K_d^{\star}$ anchor (mW\,m$^{-1}$K$^{-1}$) & \kdstarAfif & \kdstarAsev\\
best coupled RMSE (K) & \rmsecKAfif & \rmsecKAsev\\
best cp RMSE (K) & \rmsecpAfif & \rmsecpAsev\\
Hayne continuous RMSE (K) & \rmsehayAfif & 0.54\\
homogeneous RMSE (K) & \rmsehomAfif & \rmsehomAsev\\
cross-site transfer RMSE (K) & \xferatAfif & \xferatAsev\\
best structure $\rho_{1,2,3}$ (kg\,m$^{-3}$) & \bestrhocKAfif & \bestrhocKAsev\\
best boundaries $z_1/z_2$ (cm) & \bestzcKAfif & \bestzcKAsev\\
bootstrap 95\% CI, coupled (K) & \bootcicKAfif & \bootcicKAsev\\
exact-winner stability & \stabcKAfif\% & \stabcKAsev\%\\
horizon max / insolation lost & \horAfif$^\circ$ / \lossAfif\% & \horAsev$^\circ$ / \lossAsev\%\\
shadow cooling at sensors (K) & \shadowcoolAfif & \shadowcoolAsev\\
terrain IR warming (K) & $+\svfdtirAfif$ & $+\svfdtirAsev$\\
\ \ + reflected sunlight (K) & $+\svfdtallAfif$ & $+\svfdtallAsev$\\
SVF & \svfAfif & \svfAsev\\
$Q_b$ (mW\,m$^{-2}$) & 21 & 16\\
drive-core density (kg\,m$^{-3}$, Grott et al.) & ${\sim}1825$ & ${\sim}1960$\\
\bottomrule
\end{tabular}
\end{center}
\noindent Sweep: 650 configurations ($z_1\in\{5,10,20\}$ cm $\times$
$z_2\in\{30,50,80\}$ cm $\times$ 18 density triples $\times$ 2 branches);
DEM: LOLA 16 ppd (${\approx}1.9$ km/px); bootstrap: 10\,000 draws
(sensor resampling + noise); solver: ${\sim}0.6$ s per equilibrium.

% ── 6. Q&A ──────────────────────────────────────────────────────────────────
\shead{6\enspace The hard questions --- and answers that survive}
Built from a simulated five-reviewer panel; the first three are the attacks
most likely to be raised by an HFE-literate visitor.

\Q{Your coupled branch is anchored at $K_d^{\star}$, which was itself fitted
to these same sensors. Isn't the improvement circular?}
\A The anchor sets one \emph{endpoint} of the $K_c(\rho)$ line --- it does not
fix the answer. Three observations break the circularity. First, the cp
branch carries the \emph{same} anchor ($K_d=K_d^{\star}$ exactly) and still
fits far worse at A15 (\rmsecpAfif\ vs \rmsecKAfif\ K) --- so the gain comes
from the profile \emph{shape} density builds, not from re-using the anchor.
Second, the coupled winner actually \emph{moves away} from the anchor: at A17
the effective deep conductivity is $K_c(1900)=7.99$, not $7.08$ --- the data
pulled it off the calibration point. Third, the structure is validated
externally: transfer between sites and agreement with the measured drive-core
densities, neither of which the anchor can manufacture. What I concede: $K_d$
and deep density are partially degenerate from a single steady gradient
($\propto Q_b/K$); the parent study fixed density and retrieved $K_d$, this
study fixes $K_d$ and retrieves density --- conditional slices of one joint
space. The joint retrieval is the identified next step.

\Q{Your $10\times$ improvement is against a homogeneous column nobody uses.
Against Hayne, A17 improves only $1.5\times$. Isn't the baseline a straw man?}
\A Correct, and the poster carries both baselines in the table precisely so
the comparison is honest. The scoped statement: vs.\ Hayne continuous the
gains are ${\sim}2.3\times$ at A15 and ${\sim}1.5\times$ at A17. What the
strawman framing hides is actually the interesting asymmetry: Hayne already
fits A17 well (0.54 K) but \emph{fails} at A15 (\rmsehayAfif\ K under
shadowed forcing). The layered column fixes the site Hayne cannot fit, and at
the site Hayne can fit it retrieves densities that agree with the drive cores
--- two different kinds of success. The homogeneous row exists as a
floor, not as the competitor.

\Q{Five swept parameters against a dozen sensor temperatures --- overfitting?}
\A The sweep is a discrete grid (650 configurations), not a continuous
optimizer, and three controls address selection bias. (1) The cp branch is a
built-in null experiment: same parameter freedom, no conductivity channel ---
it barely moves, so the freedom alone does not produce the fit. (2)
Cross-site transfer is a held-out test: a structure tuned at one site scored
at the other. (3) The 10k-draw bootstrap resamples sensors and noise; the
improvement CIs stay below homogeneous at both sites. I quote the surviving
\emph{class} (denser surface, shallow $z_2$), not the exact triple ---
the A15 triple itself survives only \stabcKAfif\% of draws and I say so on
the poster.

\Q{If the exact A15 answer survives only 13\% of draws, what did you
actually retrieve?}
\A A class, with quantified stability: at A17 the exact configuration recurs
in \stabcKAsev\% of draws and its density triple (\bestrhocKAsev) in 88\%
(only $z_2$ wobbles between 30 and 50 cm); at A15 the deep density trades
between 1800 and 1900 while the pattern --- surface denser than Hayne's
$\rho_s$, consolidation by ${\sim}30$ cm --- recurs across draws. Point estimates in
the figure legend, class-level claims in the conclusions.

\Q{Insolation losses of 1.16\%/0.18\% are tiny. Why call topography
first-order?}
\A Because the deep sensors integrate the \emph{timing} of the loss, not its
annual average: shadowing deletes the first and last slice of every lunar
day, when the radiative $T^3$ term is weakest, so the effect rectifies
downward into \shadowcoolAfif\ K at the A15 sensors --- the same order as the
total misfit being explained. The opposing channel, terrain radiation, is
independently $+\svfdtallAfif$/$+\svfdtallAsev$ K (upper bound). At A15
cooling wins; at A17 they nearly cancel (net ${\sim}+0.3$ K). ``First-order''
means each channel is $\gtrsim$1 K-scale at A15 and must be computed --- not
that the net is always large. I concede the A17 net is small and
sign-sensitive to the terrain-temperature approximation.

\Q{Is a 1.9 km/pixel DEM adequate for horizon work near massifs?}
\A No --- and we have quantified the error rather than assumed it. Re-running
the horizon trace on SLDEM2015 crops (${\sim}59$ m/px, 512 ppd) moves A15 to
$13.6^\circ$ max / 0.54\% loss (16 ppd \emph{over}-estimated the sunrise
horizon) and A17 to $14.1^\circ$ / 0.24\% (16 ppd \emph{under}-resolved the
North Massif). The poster's numbers are the 16 ppd pipeline values and are
labeled as resolution-limited; the finer-DEM propagation through the full
sweep is the identified next step, and the pipeline takes it as a drop-in.

\Q{Why aren't the terrain-radiation terms in the fits?}
\A They were quantified as a diagnostic after the sweep, under
$T_{\rm terr}\approx T_s$ (an upper bound). Bolting them onto winners tuned
without them worsens RMSE by $+0.62/+0.31$ K --- exactly what you expect when
you perturb a tuned optimum, and the reason the honest statement is ``a
re-sweep with both channels folded in is required,'' not ``the correction is
small.''

\Q{You held $\chi$ and $Q_b$ fixed. Aren't those your biggest levers?}
\A Yes --- in the parent study's error budget they dominate. This sweep is
conditional on the certified values ($\chi=2.7$; $Q_b=21/16$ mW\,m$^{-2}$).
That is a deliberate scope: one new degree of freedom (structure) against a
frozen, certified background. Varying $\chi$, $Q_b$, and structure jointly is
the full-retrieval follow-up.

\Q{Deep sensors only --- why throw away the shallow ones?}
\A The shallow thermocouples sit in or near the borestem and are perturbed by
it (conduction down the stem, altered coupling); Langseth's own analysis
treats them separately. The deep sensors (0.8--2.4 m) are below the diurnal
skin and sample the quantity the model predicts cleanly: the mean gradient.

\Q{Why equilibrium mean profiles instead of fitting the time series?}
\A At 0.8--2.4 m the diurnal wave is dead ($e$-folding ${\sim}$10 cm); what
survives is the annual-mean state plus a small annual wave. One equilibrium
temperature per sensor from quiet windows is the robust observable; the
rectification physics that connects surface extremes to the mean gradient is
carried by the $u_{\rm rect}$ closure, not discarded.

\Q{So what? What changes if you're right?}
\A Three things. Practically: landing-site thermal assessment must ray-trace
horizons and count terrain radiation --- at 1-K scale these are not optional
corrections, and both come free from a DEM. For regolith physics: density
sensing through conductivity works --- a thermal record plus the $K_c(\rho)$
relation constrains subsurface density structure, verified here against
drive cores. For the record itself: the 50-year-old HFE data still
discriminates between modern model classes --- and will calibrate the next
generation of heat-flow probes (e.g.\ LISTER on CLPS; Nagihara's line of
work connects directly).

\Q{Does this generalize beyond Apollo?}
\A The \emph{forward} pipeline (DEM horizons, shadowed forcing, layered
solver) generalizes to any airless body with topography data. The
\emph{retrieval} needs subsurface temperature data, which exists only at
these two sites --- so: forward modeling everywhere, retrieval-ready for
future in-situ payloads. I avoid the word ``validated'' beyond A15/A17.

% ── 7. concessions & glossary ───────────────────────────────────────────────
\shead{7\enspace Concede-gracefully list}
Things you should say \emph{before} the questioner does --- each is on the
poster or in this briefing, and conceding them costs nothing:
\begin{itemize}
  \item The exact A15 triple is unstable (\stabcKAfif\%); the class is the claim.
  \item At A17, Hayne continuous is already good (0.54 K) and cp nearly ties
        coupled --- the conductivity-lever story is A15-weighted.
  \item Terrain-radiation numbers are an upper-bound diagnostic, not yet in
        the sweep.
  \item 16 ppd horizons are resolution-limited; SLDEM2015 shifts them (both
        directions) and is the planned upgrade.
  \item $K_d$ and deep density are degenerate from one gradient; joint
        retrieval pending.
  \item $\chi$ and $Q_b$ held fixed at certified values.
\end{itemize}

\shead{8\enspace Glossary (terms on the poster)}
\begin{description}\small
  \item[HFE] Apollo Heat Flow Experiment --- thermometer strings in two
    boreholes per site (A15: 1971--77; A17: 1972--77).
  \item[$K_d^{\star}$] certified site-specific deep conductivity from the
    parent retrieval (\kdstarAfif/\kdstarAsev\ mW\,m$^{-1}$K$^{-1}$).
  \item[cp branch / coupled branch] density changes heat capacity only /
    density also sets contact conductivity via $K_c(\rho)$.
  \item[SVF] sky-view factor, $\overline{\cos^2\theta_h}$ --- fraction of sky
    visible; $(1-\mathrm{SVF})$ is terrain.
  \item[ppd] pixels per degree (DEM resolution); 16 ppd $\approx$ 1.9 km/px
    at the equator.
  \item[Rectification zone / $u_{\rm rect}$] upper ${\sim}55$ cm where the
    $T^3$ nonlinearity converts oscillation into a net downward flux;
    $u_{\rm rect}$ is that flux, subtracted from $Q_b$ in the closure.
  \item[Periodic equilibrium] the repeating annual temperature cycle a
    column settles into under periodic forcing; all scoring happens there.
  \item[16/21 mW\,m$^{-2}$] $Q_b$, basal heat flux (A17/A15), Langseth's
    revised values.
\end{description}

\vspace{4pt}
{\color{teal}\rule{\linewidth}{0.9pt}}
\begin{center}\footnotesize\color{dim}
Numbers: generated macros (poster\_numbers.tex) + results archives
(aogs\_density\_study, aogs\_crossite, aogs\_winner\_bootstrap,
aogs\_svf\_terms). Full plain-language walkthrough: aogs\_study\_explainer.pdf.
\end{center}

\end{document}
